\section{Modelloversikt --- formål, forklaring og Coop-tilknytning}
\label{sec:models}
\settermlabel{terms:lec09}

\subsection*{Oversikt og lesemåte}
Denne delen samler alle de sentrale modellene fra pensum i ett oppslag.
Hver modell følger samme tre-delte struktur:

\begin{enumerate}
  \item \textbf{Formål} --- hva modellen \emph{brukes til}, og hvilken type
    spørsmål den hjelper å besvare. Dette er det viktigste å treffe på
    eksamen: feil modell på riktig spørsmål er den vanligste fellen.
  \item \textbf{Forklaring} --- de bestanddelene modellen består av, hvordan
    de henger sammen, og hvilke poenger som er kritiske å huske.
  \item \textbf{Coop-tilknytning} --- konkret kobling til Coop-casen, slik
    at modellen kan brukes som argument i prosjektforsvaret.
\end{enumerate}

\keypoint{Pointer for prosjektforsvar: bruk denne delen som verktøykasse.
Når sensor stiller et spørsmål, identifiser først hvilken modell svaret
naturlig forankres i (Formål), bruk strukturen i modellen som skjelett
(Forklaring), og avslutt alltid med Coop-spesifikk anvendelse
(Coop-tilknytning). Dette gjør svaret pensumforankret uten å bli teoretisk.}

\subsection*{Pensummodellenes plassering --- makro/meso/mikro vs.\ teknologi/forretning}

\begin{figure}[H]
  \centering
  \includegraphics[width=\linewidth]{models-overview-matrix.png}
  \caption{Oversikt over pensumets modeller plassert i matrisen med
  \emph{analyse-enhed} (makro $\rightarrow$ meso $\rightarrow$ mikro) på den
  vertikale aksen og \emph{fokus} (teknologi $\leftrightarrow$ forretning) på
  den horisontale. Bruk plasseringen som første sortering når sensor stiller
  et spørsmål: identifiser nivå og fokus, så ligger riktig modell typisk i
  den kvadranten.}
\end{figure}

% =====================================================================
\subsection{Johnson et al. --- Cultural Web / Paradigme-modellen}
\label{model:cultural-web}

\begin{center}
\includegraphics[width=0.75\linewidth]{model-cultural-web.png}
\end{center}

\subsubsection*{Formål}
Modellen brukes til å diagnostisere \emph{hvor dypt} en strategisk endring
må gripe inn. Den skiller mellom det organisasjonen \emph{sier} den tror på,
og det organisasjonen \emph{faktisk} tror på, og forklarer dermed hvorfor
strategiske endringer ofte stopper opp selv om beslutningen er tatt.

\subsubsection*{Forklaring}
Modellen er bygget i fire konsentriske lag --- jo dypere lag, jo vanskeligere
å endre:
\begin{itemize}
  \item \textbf{Values} --- de erklærte, synlige og formelle verdiene.
    Disse står i årsrapporten og på veggen.
  \item \textbf{Beliefs} --- ubevisste oppfatninger om hvordan verden
    fungerer. Disse er ikke skrevet ned, men styrer hverdagslige valg.
  \item \textbf{Behaviours} --- den faktiske atferden i praksis. Det som
    skjer når ingen ser på.
  \item \textbf{Paradigm} --- de grunnleggende antagelsene ingen stiller
    spørsmål ved. Den «selvfølgelige» logikken som hele organisasjonen
    deler.
\end{itemize}

\keypoint{Jo dypere laget er, jo vanskeligere å endre. \emph{Strategic drift}
oppstår når paradigmet ikke utfordres mens omverdenen forandrer seg.}

\subsubsection*{Coop-tilknytning}
\begin{itemize}
  \item \textbf{Values:} Coops kooperative eierskap og medlemsdemokrati ---
    erklært og synlig.
  \item \textbf{Beliefs:} «Fysiske butikker er kjernen, digitalt er støtte.»
    Ikke skrevet ned, men styrer ressursallokering.
  \item \textbf{Behaviours:} Langsomme, konsensusbaserte beslutninger på
    tvers av kjedene.
  \item \textbf{Paradigm:} «IT er drift, ikke strategi.» Dette er den
    underliggende grunnen til den teknologiske gjelden Hedman \&
    Bjørn-Andersen (2016) dokumenterer.
\end{itemize}
Den anbefalte kursen i prosjektet krever et paradigmeskift, ikke bare nye
systemer.

% =====================================================================
\subsection{PESTEL --- analyse av makromiljøet}
\label{model:pestel}

\begin{center}
\includegraphics[width=0.75\linewidth]{model-pestel.png}
\end{center}

\subsubsection*{Formål}
PESTEL brukes til å \emph{kartlegge} de eksterne makrofaktorene som kan
påvirke en virksomhet. Det er ikke en analysemodell i seg selv, men en
strukturert scanning som genererer input til Five Forces og SWOT. Poenget
er ikke å liste alt opp, men å identifisere faktorene med størst
sannsynlighet for å påvirke virksomheten.

\subsubsection*{Forklaring}
Seks dimensjoner:
\begin{itemize}
  \item \textbf{Political} --- regulering, handelspolitikk, politisk
    stabilitet.
  \item \textbf{Economic} --- vekst, inflasjon, forbrukeratferd, kjøpekraft.
  \item \textbf{Social} --- demografiske og kulturelle trender.
  \item \textbf{Technological} --- ny teknologi, digitalisering, AI.
  \item \textbf{Ecological} --- klima, bærekraft, ressursforbruk.
  \item \textbf{Legal} --- lovgivning, GDPR, konkurranserett.
\end{itemize}

\subsubsection*{Coop-tilknytning}
\begin{itemize}
  \item \textbf{Political:} Kooperativt eierskap gir spesielle
    governance-betingelser.
  \item \textbf{Economic:} Press på forbrukernes disponible inntekt
    driver etterspørsel mot discount-segmentet og 365discount.
  \item \textbf{Social:} Stigende forventning til convenience og digitale
    handleopplevelser, særlig blant yngre segmenter.
  \item \textbf{Technological:} Smartphone-penetrasjon som forutsetning
    for Coop App; AI som mulighet for å utnytte førstepartsdata.
  \item \textbf{Ecological:} Stigende forventning til bærekraftige
    produkter og kooperativets ansvarsprofil.
  \item \textbf{Legal:} GDPR er direkte relevant for hvordan Coop kan
    samle inn og bruke atferdsdata i plattformen.
\end{itemize}
\keypoint{Technological og Legal er de mest direkte relevante
dimensjonene for Coop-casen.}

% =====================================================================
\subsection{Porters Five Forces}
\label{model:five-forces}

\begin{center}
\includegraphics[width=0.75\linewidth]{model-five-forces.png}
\end{center}

\subsubsection*{Formål}
Brukes til å vurdere \emph{hvor attraktiv en industri er}, og hvor
konkurransepresset kommer fra. Analyseenheten er \emph{industrien}, ikke
virksomheten.

\subsubsection*{Forklaring}
Fem konkurransekrefter:
\begin{itemize}
  \item \textbf{Rivalisering i industrien} --- intensiteten av
    konkurransen mellom eksisterende aktører.
  \item \textbf{Trusselen fra nye inntrengere} --- hvor vanskelig det er
    å komme inn på markedet (adgangsbarrierer).
  \item \textbf{Substituerende produkter} --- alternativer som løser
    samme behov på en annen måte.
  \item \textbf{Leverandørenes forhandlingskraft} --- avhengighet av
    leverandører.
  \item \textbf{Kundenes forhandlingskraft} --- hvor lett kundene kan
    skifte.
\end{itemize}

\keypoint{Porter tilføyde selv en sjette kraft i 1990 ---
\emph{komplementorer} --- aktører som tilfører verdi til plattformen.}

\subsubsection*{Coop-tilknytning}
\begin{itemize}
  \item \textbf{Rivalisering:} Høy. Salling Group investerer massivt i
    digitale løsninger og har modernisert supply chain hurtigere enn Coop.
  \item \textbf{Nye inntrengere:} Moderat på fysisk detaljhandel
    (høye adgangsbarrierer), men lav på digital side hvor
    plattformsaktører som Amazon Fresh potensielt kan entre uten butikker.
  \item \textbf{Substitutter:} Måltidskasser (HelloFresh), takeaway og
    direkte leverandørplattformer.
  \item \textbf{Leverandører:} Moderat. Store brands har forhandlingskraft,
    men Coops 1,8 mio. brukere gir motvekt.
  \item \textbf{Kunder:} Høy. Ingen kontrakt, lave switching costs,
    Salling-appen er ett sveip unna.
  \item \textbf{Komplementorer:} Brands og leverandører i Coop App er
    nettopp komplementorer --- ikke bare leverandører. Nyttig nyansering
    i plattformsanalyse-kontekst.
\end{itemize}

% =====================================================================
\subsection{Industriers livssyklus og De Fem Kreftene}
\label{model:industry-lifecycle}

\begin{center}
\includegraphics[width=0.8\linewidth]{model-industry-lifecycle.png}
\end{center}

\subsubsection*{Formål}
Brukes til å vise at \emph{Five Forces ser forskjellig ut} avhengig av
hvor i livssyklusen en industri befinner seg. Konkurransedynamikken
endrer seg over tid, og strategiske valg må tilpasses fasen.

\subsubsection*{Forklaring}
Fem faser:
\begin{itemize}
  \item \textbf{Development} --- lav rivalisering, høy differensiering,
    innovasjon avgjørende. Få spillere, markedet defineres ennå.
  \item \textbf{Growth} --- fortsatt lav rivalisering, mange nye
    inntrengere, vekst avgjørende. Alle vil ha en bit.
  \item \textbf{Shake-out} --- rivalisering stiger, svake aktører
    forlater markedet, finansiell og ledelsesmessig styrke avgjørende.
  \item \textbf{Maturity} --- sterkere kunder, standardiserte produkter,
    høye adgangsbarrierer, markedsandel og kostnader avgjørende.
  \item \textbf{Decline} --- ekstrem rivalisering, mange exits,
    priskonkurranse dominerer. Kostnadsstyring og commitment teller.
\end{itemize}

\subsubsection*{Coop-tilknytning}
Dansk dagligvarehandel \emph{fysisk} er i \textbf{Maturity}. Det forklarer
det massive konkurransepresset fra Salling Group, de lave marginene og
omsetningsfallet. Den \emph{digitale} plattformsdelen av industrien er
derimot i \textbf{Growth} eller tidlig \textbf{Shake-out} --- det er
nettopp dette som gjør det til riktig tidspunkt for Coop å konsolidere
sin digitale posisjon framfor å trekke seg.

% =====================================================================
\subsection{VRIO --- ressursbasert analyse}
\label{model:vrio}

\begin{center}
\includegraphics[width=0.7\linewidth]{model-vrio-1.png}\\[6pt]
\includegraphics[width=0.7\linewidth]{model-vrio-2.png}
\end{center}

\subsubsection*{Formål}
VRIO brukes til å vurdere \emph{én ressurs eller kapabilitet av gangen}
og avgjøre om den gir varig konkurransefordel. Modellen er internt
fokusert (motstykket til Five Forces' eksterne fokus).

\subsubsection*{Forklaring}
Fire spørsmål stilles sekvensielt:
\begin{itemize}
  \item \textbf{Valuable} --- bidrar ressursen til å imøtekomme en
    mulighet eller avverge en trussel?
  \item \textbf{Rare} --- besitter få eller ingen konkurrenter samme
    ressurs?
  \item \textbf{Inimitable} --- er den vanskelig og kostbar å kopiere
    eller substituere?
  \item \textbf{Organized} --- er organisasjonen rigget til å utnytte
    ressursen strategisk?
\end{itemize}

\subsubsection*{Coop-tilknytning}
Coops førstepartsdata fra 1,8 mio. brukere scorer sannsynligvis \emph{ja}
på V, R og I. Dataen er verdifull, sjelden og vanskelig å kopiere. Men
\textbf{O} er det kritiske spørsmålet. Er Coop organisert til å utnytte
dataen strategisk? Rapporten antyder \emph{nei}, fordi provider-siden er
svak og governance er underutviklet. Det plasserer Coop i
\textbf{«uutnyttet konkurransefordel»} --- og det er nettopp dette
anbefalingene forsøker å endre.

% =====================================================================
\subsection{SWOT}
\label{model:swot}

\begin{center}
\includegraphics[width=0.7\linewidth]{model-swot.png}
\end{center}

\subsubsection*{Formål}
SWOT samler intern og ekstern analyse i én oversikt. Brukes som
\emph{oppsummering}, ikke som startpunkt. SWOT er ikke en analysemodell i
seg selv, men en sammenstilling av VRIO (intern) og Five Forces/PESTEL
(ekstern).

\subsubsection*{Forklaring}
Fire kategorier:
\begin{itemize}
  \item \textbf{Strengths} --- interne styrker.
  \item \textbf{Weaknesses} --- interne svakheter.
  \item \textbf{Opportunities} --- eksterne muligheter.
  \item \textbf{Threats} --- eksterne trusler.
\end{itemize}

\subsubsection*{Coop-tilknytning}
\begin{itemize}
  \item \textbf{Strengths:} 1,8 mio. brukere, førstepartsdata, Lobycos
    OEM-plattform, sterk brand-gjenkjenning.
  \item \textbf{Weaknesses:} Teknologisk gjeld, svak provider-side,
    lukket partnerstruktur, underskudd på 232 mio. DKK, lukket governance.
  \item \textbf{Opportunities:} Voksende retail media-marked, AI-drevet
    personalisering, Lobycos internasjonale OEM-ekspansjon til SPAR UK
    og SPAR Hungary.
  \item \textbf{Threats:} Salling Groups digitale investeringer,
    «appflation», potensiell plattformdisrupsjon fra Amazon Fresh eller
    lignende, GDPR-regulering.
\end{itemize}
\keypoint{Sensor blir ikke imponert av en SWOT alene, men av en SWOT som
er forankret i de underliggende analysene.}

% =====================================================================
\subsection{TOWS --- fra analyse til handling}
\label{model:tows}

\begin{center}
\includegraphics[width=0.7\linewidth]{model-tows.png}
\end{center}

\subsubsection*{Formål}
TOWS er \emph{SWOT vendt om til handling}. Hvor SWOT beskriver
situasjonen, genererer TOWS strategiske muligheter ved å krysse de fire
feltene parvis.

\subsubsection*{Forklaring}
Fire kombinasjoner:
\begin{itemize}
  \item \textbf{SO --- Strengths + Opportunities:} bruk styrker til å
    utnytte muligheter.
  \item \textbf{ST --- Strengths + Threats:} bruk styrker til å avverge
    trusler.
  \item \textbf{WO --- Weaknesses + Opportunities:} overvinn svakheter
    for å utnytte muligheter.
  \item \textbf{WT --- Weaknesses + Threats:} minimer svakheter for å
    unngå trusler.
\end{itemize}

\subsubsection*{Coop-tilknytning}
\begin{itemize}
  \item \textbf{SO:} Brukerdata + voksende retail media-marked =
    profesjonaliser retail media-modellen $\to$ \emph{Anbefaling 3}.
  \item \textbf{ST:} Samlet plattform + appflation-trussel = fasthold
    én app framfor kjedespesifikke løsninger $\to$ \emph{Anbefaling 2}.
  \item \textbf{WO:} Svak provider-side + retail media-marked =
    strukturer provider-adgang via kontrollert system.
  \item \textbf{WT:} Teknologisk gjeld + plattformdisrupsjon = fasthold
    investeringene i Lobyco og Coop App $\to$ \emph{Anbefaling 1}.
\end{itemize}
\keypoint{TOWS er særlig sterk fordi den viser at anbefalingene
\emph{ikke er tilfeldige}. Alle tre anbefalinger kan forankres direkte i
TOWS-logikken, selv om vi ikke brukte modellen eksplisitt i rapporten.}

% =====================================================================
\subsection{Verdikjede og verdikjedesystemer (Porter)}
\label{model:value-chain}

\begin{center}
\includegraphics[width=0.75\linewidth]{model-value-chain-1.png}\\[6pt]
\includegraphics[width=0.75\linewidth]{model-value-chain-2.png}
\end{center}

\subsubsection*{Formål}
Verdikjeden brukes til å bryte virksomheten ned i \emph{aktiviteter} som
skaper kundeverdi, og identifisere hvor konkurransefordel oppstår.
Verdikjedesystemet utvider perspektivet til å inkludere leverandørenes og
kundenes verdikjeder --- og åpner for at konkurransefordel kan oppstå i
\emph{grenseflaten} mellom kjedene.

\subsubsection*{Forklaring}
\textbf{Primære aktiviteter} skaper direkte verdi:
\begin{itemize}
  \item \textbf{Inbound logistics} --- mottak og lagring.
  \item \textbf{Operations} --- omdanning til ferdig produkt.
  \item \textbf{Outbound logistics} --- levering til kunden.
  \item \textbf{Marketing and sales} --- kommunikasjon og salg.
  \item \textbf{Service} --- ettersalg og kundeservice.
\end{itemize}

\textbf{Støtteaktiviteter} understøtter alle primære:
\begin{itemize}
  \item \textbf{Firm infrastructure} --- ledelse, økonomi, governance.
  \item \textbf{Human resource management} --- rekruttering, kompetanse.
  \item \textbf{Technology development} --- F\&U, IT.
  \item \textbf{Procurement} --- innkjøp.
\end{itemize}

\subsubsection*{Coop-tilknytning}
\begin{itemize}
  \item \textbf{Inbound:} Distribusjonssentre og vareflyt.
  \item \textbf{Operations:} Butikkdrift.
  \item \textbf{Outbound:} Fysiske butikker som primær kanal ---
    coop.dk ble lukket i 2023.
  \item \textbf{Marketing/sales:} Coop App, bonusprogram, personaliserte
    tilbud.
  \item \textbf{Service:} Scan \& Pay, handlelister, medlemsfordeler.
  \item \textbf{Technology development:} Lobyco og Coop App er Coops
    primære teknologiinvestering.
\end{itemize}
Coop App opererer nettopp i \emph{grenseflaten} mellom leverandørenes og
kundenes verdikjeder --- det er det som gjør den til en plattform framfor
en intern IT-løsning. Dette er også grunnen til at verdikjedeanalyse
alene ikke er tilstrekkelig for Coop App; Eisapes Canvas trengs for å
håndtere de flere verdikjedene samtidig.

% =====================================================================
\subsection{Five Building Blocks of Digital Transformation (Ross et al.)}
\label{model:five-building-blocks}

\begin{center}
\includegraphics[width=0.75\linewidth]{model-five-building-blocks.png}
\end{center}

\subsubsection*{Formål}
Modellen brukes til å diagnostisere \emph{hvilke fundamentale
byggeklosser} som må være på plass for at digital transformasjon skal
lykkes. Den forklarer hvorfor mange transformasjonsprosjekter feiler ---
fordi én eller flere blokker mangler.

\subsubsection*{Forklaring}
Fem byggeklosser:
\begin{itemize}
  \item \textbf{Operational Backbone} --- integrerte bakgrunnssystemer,
    ERP, supply chain. Fundamentet alt annet hviler på.
  \item \textbf{Digital Platform} --- gjenbrukbare teknologikomponenter
    interne team kan bygge på.
  \item \textbf{External Developer Platform} --- åpner infrastrukturen
    for tredjepartsutviklere.
  \item \textbf{Shared Customer Insights} --- delt kundekunnskap på
    tvers av forretningsenheter.
  \item \textbf{Accountability Framework} --- klart eierskap og ansvar
    for digitale produkter.
\end{itemize}

\subsubsection*{Coop-tilknytning}
\begin{itemize}
  \item \textbf{Operational Backbone:} Fragmentert og foreldet i 2016
    ifølge Hedman \& Bjørn-Andersen.
  \item \textbf{Digital Platform:} Lobyco er forsøket på dette.
  \item \textbf{External Developer Platform:} Ikke her ennå ---
    partnerstrukturen er lukket og bilateral.
  \item \textbf{Shared Customer Insights:} Potensialet finnes med 1,8
    mio. brukeres data, men uklart om det deles effektivt på tvers av
    kjedene.
  \item \textbf{Accountability Framework:} Uklart pga. todelt eierskap
    mellom Lobyco og Coop Danmark.
\end{itemize}
\keypoint{Coops transformasjon er \emph{gradvis og uferdig, ikke
mislykket}. Retail media-anbefalingen forutsetter at Shared Customer
Insights og Accountability Framework styrkes. God modell for å vise
\emph{hvorfor} implementering er vanskelig hos Coop.}

% =====================================================================
\subsection{CVP --- Johnson et al.s forretningsmodellramme}
\label{model:cvp}

\begin{center}
\includegraphics[width=0.75\linewidth]{model-cvp.png}
\end{center}

\subsubsection*{Formål}
Brukes til å analysere en forretningsmodell som \emph{fire sammenhengende
elementer} som må være konsistente med hverandre. CVP er utgangspunktet;
de tre andre må understøtte den. Modellen forklarer hvorfor en sterk
verdiproposisjon \emph{alene} ikke er nok (jf.\ Tata Nano-eksempelet).

\subsubsection*{Forklaring}
Fire elementer:
\begin{itemize}
  \item \textbf{Customer Value Proposition} --- target customer, job to
    be done, offering.
  \item \textbf{Profit Formula} --- revenue model, cost structure,
    margin model, resource velocity.
  \item \textbf{Key Resources} --- people, technology, information, brand.
  \item \textbf{Key Processes} --- prosesser, regler og metrikker, normer.
\end{itemize}

\subsubsection*{Coop-tilknytning}
\begin{itemize}
  \item \textbf{CVP:} Target = 1,8 mio. app-brukere. Job = handle
    effektivt, spare penger, få personaliserte opplevelser. Offering =
    Scan \& Pay, personaliserte tilbud, bonusprogram, handlelister.
  \item \textbf{Profit Formula:} Retail media-inntekter + indirekte
    omsetning via økt butikkbesøksfrekvens (app-brukere handler 50\%
    hyppigere). Lav direkte margin på app, høy strategisk verdi via
    datadrevet butikkatferd.
  \item \textbf{Key Resources:} Lobycos utviklingsorganisasjon,
    Coop App-infrastruktur, OEM-plattform, atferdsdata fra 1,8 mio.
    brukere (mest kritisk), Coops kooperative identitet.
  \item \textbf{Key Processes:} Plattformsutvikling, kampanjeavvikling,
    dataanalyse, partnerstyring; KPI-er for app-engagement, retail
    media-omsetning, butikkbesøksfrekvens.
\end{itemize}
\keypoint{Coops svakhet er at \emph{Profit Formula} på retail media-siden
ennå ikke er profesjonalisert, selv om CVP og Key Resources er på plass.
Det er nøyaktig dette Anbefaling 3 adresserer.}

% =====================================================================
\subsection{Business Model Canvas (Osterwalder)}
\label{model:bmc}

\begin{center}
\includegraphics[width=0.85\linewidth]{model-bmc-2.png}\\[6pt]
\includegraphics[width=0.6\linewidth]{model-bmc-1.png}
\end{center}

\subsubsection*{Formål}
BMC er et visuelt verktøy til å beskrive en virksomhets forretningsmodell
i ni byggesteiner. Designet for \emph{pipeline-virksomheter} som skaper
og leverer verdi i én retning fra produsent til kunde.

\subsubsection*{Forklaring}
Ni byggesteiner: Customer Segments, Value Propositions, Channels,
Customer Relationships, Revenue Streams, Key Resources, Key Activities,
Key Partnerships, Cost Structure.

\subsubsection*{Coop-tilknytning}
BMC fungerer fint for Coops fysiske butikkdrift, men er \emph{utilstrekkelig}
for Coop App. Plattformslogikken med fire aktørgrupper (Owner, Consumer,
Provider, Partner) med motstridende interesser fanges ikke i den lineære
pipeline-strukturen. Det er nettopp derfor en plattform-spesifikk Canvas
(Eisape) er nødvendig.

% =====================================================================
\subsection{Weill \& Woerner (2018) --- Digital Business Model Framework}
\label{model:weill-woerner}

\begin{center}
\includegraphics[width=0.75\linewidth]{model-weill-woerner.jpeg}
\end{center}

\subsubsection*{Formål}
DBM brukes til å \emph{plassere} en virksomhet i et todimensjonalt
landskap basert på forretningsdesign og kundekunnskap, og deretter
identifisere \emph{retning} (hvor virksomheten beveger seg).

\subsubsection*{Forklaring}
To akser danner fire kvadranter:
\begin{itemize}
  \item \textbf{X-akse:} verdikjedebasert $\to$ økosystembasert.
  \item \textbf{Y-akse:} partiell $\to$ komplett kundekunnskap.
\end{itemize}
Fire posisjoner:
\begin{itemize}
  \item \textbf{Supplier} --- partiell kunnskap, verdikjede. Selger via
    andres kanaler.
  \item \textbf{Omnichannel} --- komplett kunnskap, verdikjede. Eier
    kunderelasjonen digitalt, men opererer fortsatt lineært.
  \item \textbf{Modular Producer} --- partiell kunnskap, økosystem.
    Leverer komponenter til andres plattformer.
  \item \textbf{Ecosystem Driver} --- komplett kunnskap, økosystem.
    Faciliterer interaksjoner mellom mange aktører.
\end{itemize}

\subsubsection*{Coop-tilknytning}
\begin{itemize}
  \item \textbf{Nåværende posisjon:} \emph{Omnichannel}. Coop eier
    kunderelasjonen digitalt via appen og har komplett kundekunnskap via
    førstepartsdata, men opererer fortsatt i verdikjedelogikk med
    fysiske butikker som primær kanal.
  \item \textbf{Risikoscenario:} \emph{Supplier}. Hvis app-investeringene
    reduseres og kundekunnskapen mistes til tredjeparter, glir Coop mot
    å levere varer via andres plattformer.
  \item \textbf{Ambisjonsscenario:} \emph{Ecosystem Driver}. Lobycos
    OEM-ekspansjon til SPAR UK og SPAR Hungary peker i denne retningen,
    men krever åpent økosystembasert design og investeringskapasitet
    som er uforenlig med underskuddet på 232 mio. DKK.
  \item \textbf{Modular Producer:} Lobyco beveger seg mot dette i sin
    OEM-rolle overfor internasjonale partnere.
\end{itemize}

% =====================================================================
\subsection{Network effects (Nettverkseffekter)}
\label{model:network-effects}

\begin{center}
\includegraphics[width=0.8\linewidth]{model-network-effects.png}
\end{center}

\subsubsection*{Formål}
Nettverkseffekter brukes til å forklare \emph{hvorfor plattformer blir
strategiske giganter}, og hvorfor «winner takes most»-dynamikker oppstår
i plattformsmarkeder.

\begin{definition}{Network effects}
Nettverkseffekter refererer til den innvirkning antallet brukere på en
plattform har på verdien som skapes for hver bruker.
\end{definition}

\subsubsection*{Forklaring}
Fire varianter:
\begin{itemize}
  \item \textbf{Positiv same-side:} Flere brukere av samme type øker
    verdien for hverandre. Eks.: Instagram --- flere brukere gir mer
    innhold og bedre opplevelse.
  \item \textbf{Positiv cross-side:} Flere brukere på den ene siden
    tiltrekker den andre. Eks.: Uber --- flere sjåfører gir kortere
    ventetid, tiltrekker flere passasjerer.
  \item \textbf{Negativ same-side:} For mange av samme type skaper
    konkurranse og reduserer verdi. Eks.: Tinder --- mange av ett
    kjønn øker konkurransen og senker matchraten.
  \item \textbf{Negativ cross-side:} For mange på én side forringer
    opplevelsen for den andre. Eks.: DBA --- for mange selgere med
    lavkvalitetsvarer forverrer kjøperopplevelsen.
\end{itemize}

\subsubsection*{Coop-tilknytning}
Coop App utnytter primært \emph{positive cross-side effekter}. Flere
brands og leverandører på provider-siden gir bedre tilbud til
consumer-siden, som tiltrekker flere brukere, som igjen gir providers
tilgang til en større base. Dette er den selvforsterkende mekanismen som
gjør plattformen verdifull, og det er nøyaktig denne mekanismen som
forsvinner hvis Coop går ned til Supplier-posisjon i DBM.

% =====================================================================
\subsection{Allwein et al.\ (2020) --- Platform Business Model Canvas}
\label{model:allwein}

\begin{center}
\includegraphics[width=0.8\linewidth]{model-allwein.jpeg}
\end{center}

\subsubsection*{Formål}
Allwein-canvas brukes til å analysere \emph{én primær interaksjon} mellom
to segmenter på en plattform. Designet for \emph{enkle two-sided}
plattformer.

\subsubsection*{Forklaring}
Enkel og symmetrisk struktur. Fokuserer på kjerneinteraksjonen mellom to
sider. Mindre fokus på governance og kompleksitet. Passer godt til
Airbnb-, Tinder- eller eBay-lignende plattformer hvor det reelt er to
sider som møtes.

\subsubsection*{Coop-tilknytning}
\emph{Ikke egnet} for Coop App. Coop App er ikke en simpel two-sided
plattform: det er fire sideordnede aktørgrupper (Owner, Consumer,
Provider, Partner) med motstridende interesser. Allwein-canvas ville
tvinge oss til å redusere kompleksiteten kunstig og miste den
strategiske diagnosens presisjon. Brukes derfor som
\emph{kontrastmodell} i prosjektet --- for å vise \emph{hvorfor} Eisape
er valgt.

% =====================================================================
\subsection{Lubinski et al.\ (2020) --- Platform Business Model Canvas}
\label{model:lubinski}

\begin{center}
\includegraphics[width=0.8\linewidth]{model-lubinski.png}
\end{center}

\subsubsection*{Formål}
Lubinski-canvas brukes til å analysere \emph{vekst, governance og
utvikling over tid} for en plattform. Velegnet når plattformen er i
aktiv vekst- og reguleringsmodning.

\subsubsection*{Forklaring}
Inndelt i tre lag:
\begin{itemize}
  \item \textbf{Core} --- selve plattformens kjernemekanikk.
  \item \textbf{Business} --- forretningsmodellen som omslutter kjernen.
  \item \textbf{Institutional setup} --- governance, kvalitetsstandarder,
    regulering.
\end{itemize}
Sterkt fokus på utviklingen fra åpent mot kontrollert system.

\subsubsection*{Coop-tilknytning}
\emph{Delvis relevant.} Lubinski ville vært egnet hvis analysen
fokuserte på Lobycos vekststrategi og governance-modning over tid. Men
oppgaven er en \emph{statisk strategisk diagnose} --- ikke en
modningsanalyse. Brukes derfor som kontrastmodell, ikke som hovedverktøy.

% =====================================================================
\subsection{Davis Eisape --- Platform Business Model Canvas}
\label{model:eisape}

\begin{center}
\includegraphics[width=0.85\linewidth]{model-eisape.png}
\end{center}

\subsubsection*{Formål}
Eisape-canvas brukes til å analysere \emph{komplekse plattformer med
flere aktørgrupper} hvor governance og kuratering er sentralt. Modellen
håndterer Owner, Consumer, Provider og Partner som \emph{analytisk
adskilte} segmenter.

\subsubsection*{Forklaring}
Strukturerer plattformen rundt fire sideordnede aktørgrupper med
hver sine value propositions, value capture-mekanismer og relasjoner
til plattformeieren. Eksplisitt fokus på \emph{kuratering} ---
plattformeierens rolle i å balansere de fire gruppenes interesser.

\subsubsection*{Coop-tilknytning}
\emph{Den valgte hovedmodellen i prosjektet.} Eisape er den eneste av de
tre Platform Canvas-modellene som håndterer Coop Apps faktiske
kompleksitet:
\begin{itemize}
  \item \textbf{Owner:} Coop Danmark + Lobyco (todelt eierskap er
    kjernen i governance-utfordringen).
  \item \textbf{Consumer:} 1,8 mio. app-brukere.
  \item \textbf{Provider:} Brands og leverandører på retail media-siden.
  \item \textbf{Partner:} Shortcut, Playable og andre teknologipartnere.
\end{itemize}
Eisape gjør det mulig å vise gapet mellom \emph{skapt} og \emph{fanget}
verdi som er hovedargumentet i prosjektet.

% =====================================================================
\subsection{Wessel et al.\ --- Digital Transformation vs IT-Enabled Transformation}
\label{model:wessel}

\begin{center}
\includegraphics[width=0.8\linewidth]{model-wessel-1.png}\\[6pt]
\includegraphics[width=0.65\linewidth]{model-wessel-2.png}
\end{center}

\subsubsection*{Formål}
Modellen brukes til å \emph{skille} mellom to fundamentalt forskjellige
typer transformasjon: ekte digital transformasjon (DT) som omdefinerer
verdiproposisjonen, vs IT-Enabled Organizational Transformation (ITOT)
som understøtter den eksisterende. Sentral for diagnose: er Coop
egentlig i DT eller ITOT?

\subsubsection*{Forklaring}
Modellen har fem komponenter:
\begin{itemize}
  \item \textbf{Technological Change} --- det ytre presset:
    \emph{Environmental context} (konkurrenter, markedsforventninger) +
    \emph{Organizational context} (interne ressurser, kultur).
  \item \textbf{Transformation Agenda} --- det indre svaret:
    \emph{Existing Organizational Identity} (hvem er vi i dag) +
    \emph{New Identity Claim} (hvem vil vi bli).
  \item \textbf{Transformation Activity} --- selve transformasjonen.
    Digital teknologi enten \emph{omdefinerer} (DT) eller
    \emph{understøtter} (ITOT) verdiproposisjonen.
  \item \textbf{Imposition \& Reconciliation} --- spenningen mellom ny
    og gammel identitet. \emph{Work Practice Change} (faktisk endring i
    praksis) + \emph{Reconciliation Action} (ledelseshandlinger).
  \item \textbf{To utfall:} \emph{New Organizational Identity}
    (transformasjon lyktes) eller \emph{Reinforced Organizational
    Identity} (eksisterende identitet styrket men ikke transformert).
\end{itemize}

\subsubsection*{Coop-tilknytning}
\begin{itemize}
  \item \textbf{Environmental:} Salling Groups digitale investeringer og
    stigende forbrukerforventninger.
  \item \textbf{Organizational:} Teknologisk gjeld og kooperativ
    beslutningskultur.
  \item \textbf{Existing identity:} Kooperativ dagligvarekonsern med
    fysiske butikker som primær kanal.
  \item \textbf{New identity claim:} Datadrevet detaljhandler som
    skaper verdi via digital kundeinnsikt og retail media.
  \item \textbf{Transformation activity:} Coop App og Lobyco
    \emph{understøtter} eksisterende verdiproposisjon, men omdefinerer
    den ikke fullt ut ennå.
  \item \textbf{Reconciliation:} Konsolidering av Omnichannel-posisjonen
    er nettopp en reconciliation action.
  \item \textbf{Sannsynlig utfall:} Reinforced Organizational Identity
    på kort sikt.
\end{itemize}
\keypoint{Coop er i ITOT, ikke ekte DT --- og det er en bevisst, riktig
strategi gitt EBIT $-$215 mio. DKK og lukkingen av coop.dk.}

% =====================================================================
\subsection{Øyemodellen --- typer disruption}
\label{model:eye}

\begin{center}
\includegraphics[width=0.75\linewidth]{model-eye.png}
\end{center}

\subsubsection*{Formål}
Brukes til å \emph{identifisere og presisere} eksterne disruptive
trusler mot en incumbent. Forhindrer at «disruption» brukes som løst
buzzword --- tvinger til å spesifisere både \emph{type} og \emph{målmarked}.

\subsubsection*{Forklaring}
To akser:

\textbf{Type (venstre):}
\begin{itemize}
  \item \emph{Disruptive Technology Innovation} --- ny teknologi er
    driveren.
  \item \emph{Disruptive Radical Innovation} --- teknologi + ny
    tenkning kombinert.
  \item \emph{Disruptive Business Model Innovation} ---
    forretningsmodellen er det disruptive.
\end{itemize}

\textbf{Marked (høyre):}
\begin{itemize}
  \item \emph{Low-End Disruption} --- starter hos minst krevende kunder,
    billigere og dårligere.
  \item \emph{High-End Disruption} --- starter på toppen, beveger seg
    nedover.
  \item \emph{New Market Disruption} --- skaper et marked for
    ikke-kunder.
\end{itemize}

\subsubsection*{Coop-tilknytning}
Trusselen er primært \textbf{Disruptive Business Model Innovation +
New Market Disruption}. En plattformsaktør som Amazon Fresh kan
integrere dagligvarer i et bredere livsstilsøkosystem og nå forbrukere
Coop ikke har kontakt med i dag. Det er ikke teknologien som er
disruptiv, men forretningsmodellen som gjør dagligvarer til en
sideydelse i et større plattformsøkosystem.

Coops forsvar er den datamessige posisjonen og kunderelasjonen via
Coop App. Mistes den, mistes det eneste som kan motstå disrupsjon utenfra.

% =====================================================================
\subsection{Disruption-modellen (tre sirkler)}
\label{model:disruption-circles}

\begin{center}
\includegraphics[width=0.7\linewidth]{model-disruption-circles.png}
\end{center}

\subsubsection*{Formål}
Brukes til å \emph{presisere hva digital disruption faktisk er}, og
skille det fra beslektede begreper som ofte forveksles (digital
innovasjon, disruptiv innovasjon).

\subsubsection*{Forklaring}
Tre sirkler:
\begin{itemize}
  \item \textbf{Digital Innovation (DgI)} --- bred blå sirkel. All
    nyskapning som involverer digital teknologi. Ikke nødvendigvis
    disruptiv.
  \item \textbf{Disruptive Innovation (DI)} --- bred oransje sirkel.
    Innovasjoner som endrer et marked fundamentalt. Ikke nødvendigvis
    digital.
  \item \textbf{Digital Disruption (DgD)} --- det grønne fellesfeltet.
    Kun når innovasjonen er \emph{både digital og disruptiv samtidig}.
\end{itemize}

\keypoint{De fleste digitale innovasjoner er ikke disruptive. Mange
disruptive innovasjoner er ikke digitale. Digital disruption er det
smale overlappet hvor begge er til stede.}

\subsubsection*{Coop-tilknytning}
Coop App er \emph{digital innovasjon} men ikke \emph{digital disruption},
fordi den ikke omdefinerer paradigmet i dagligvarehandelen. Den
understøtter eksisterende forretningslogikk.

Det som potensielt \emph{truer} Coop er derimot digital disruption
\emph{utenfra} --- f.eks.\ Amazon Fresh --- som kombinerer digital
teknologi med en fundamentalt annerledes forretningsmodell som
omdefinerer hva dagligvarehandel \emph{er}.

% =====================================================================
\subsection{Digital Innovation Strategy + tre kjerneelementer}
\label{model:digital-innovation-strategy}

\begin{center}
\includegraphics[width=0.75\linewidth]{model-digital-innovation-strategy.png}
\end{center}

\subsubsection*{Formål}
Brukes til å \emph{kartlegge hvor} en virksomhets digitale
innovasjonsinnsats er plassert, og identifisere \emph{blinde flekker}
i innovasjonsstrategien.

\subsubsection*{Forklaring}
Fire felter (matrise: Invent NEW / Enhance EXISTING $\times$ Internal /
External):
\begin{itemize}
  \item \textbf{Invent NEW + Internal:} nye forretningsprosesser og
    medarbeideropplevelser.
  \item \textbf{Invent NEW + External:} nye forretningsmodeller og
    kundeopplevelser.
  \item \textbf{Enhance EXISTING + Internal:} forbedre eksisterende
    prosesser. Effektivisering.
  \item \textbf{Enhance EXISTING + External:} forbedre eksisterende
    forretningsmodell og kundeopplevelser. Iterativ utvikling.
\end{itemize}

\textbf{Tre kjerneelementer i digital innovasjon:}
\begin{itemize}
  \item \textbf{Digital technology} --- teknologien er middelet, ikke
    målet.
  \item \textbf{Recombination} --- digital innovasjon oppstår ved å
    kombinere eksisterende digitale og fysiske komponenter på nye måter.
  \item \textbf{Value Co-Creation} --- verdien skapes i samspillet
    mellom aktører, ikke av én virksomhet alene.
\end{itemize}

\subsubsection*{Coop-tilknytning}
\begin{itemize}
  \item \textbf{Enhance EXISTING + External:} Coop App. Forbedrer den
    eksisterende kundeopplevelsen via personalisering og convenience,
    men oppfinner ikke en ny forretningsmodell.
  \item \textbf{Invent NEW + External:} Lobycos OEM-ekspansjon til
    SPAR UK og SPAR Hungary --- en potensielt ny forretningsmodell
    basert på plattformslisensiering.
  \item \textbf{Retail media-anbefalingen:} en bevegelse fra Enhance
    mot Invent på den eksterne aksen, fordi den skaper en ny
    inntektsstrøm som ikke eksisterte før.
\end{itemize}
\keypoint{Coops innovasjonsinnsats er konsentrert i ett felt. En mer
ambisiøs digital strategi ville krevd innovasjon på tvers av alle fire
felter --- særlig Invent NEW internt, hvor medarbeider- og
prosessutvikling er nødvendig for å realisere anbefalingene.}

% =====================================================================
\subsection{Tawse \& Tabesh --- Strategy Implementation Framework}
\label{model:tawse-tabesh}

\begin{center}
\includegraphics[width=0.8\linewidth]{model-tawse-tabesh.png}
\end{center}

\subsubsection*{Formål}
Brukes til å \emph{analysere og forklare} hvorfor strategiimplementering
lykkes eller feiler. Identifiserer de konkrete ledelseshandlinger og
forutsetninger som må være på plass for at endring skal manifestere seg
som faktisk atferd. Målet er \emph{endret atferd} i organisasjonen ---
ikke bare en vedtatt strategi.

\subsubsection*{Forklaring}
Tre lag, ovenfra og ned:

\textbf{Lag 1 --- Conditions} (organisatoriske forutsetninger):
\begin{itemize}
  \item \textbf{Kapacitet} --- har medarbeiderne nødvendige kompetanser
    og ressurser?
  \item \textbf{Commitment} --- er organisasjonsmedlemmene motiverte og
    engasjerte?
  \item \textbf{Koordinasjon} --- er tverrgående prosesser og
    kommunikasjon på plass?
\end{itemize}

\textbf{Lag 2 --- Managerial Actions} (konkrete ledelseshandlinger):
\begin{itemize}
  \item \textbf{Strukturelle} --- belønningssystemer, prosedyrer,
    kontroll, formelle prosesser.
  \item \textbf{Interpersonelle} --- kommunikasjon, tillitsbygging,
    konsensus om strategiske mål.
\end{itemize}

\textbf{Lag 3 --- Dynamic Managerial Capabilities} (underliggende
ledelseskompetanser):
\begin{itemize}
  \item \textbf{Managerial cognition} --- forståelse for konsekvensene
    av ens valg.
  \item \textbf{Managerial social capital} --- relasjoner som gir tilgang
    til ressurser og informasjon.
  \item \textbf{Managerial human capital} --- erfaring, ekspertise,
    planleggingskompetanse.
\end{itemize}

\subsubsection*{Coop-tilknytning}
\begin{itemize}
  \item \textbf{Kapacitet:} Coop mangler sannsynligvis retail
    media-kompetanse til å eksekvere Anbefaling 3.
  \item \textbf{Commitment:} Kjede-autonomien gjør det vanskelig å sikre
    engasjement på tvers av organisasjonen.
  \item \textbf{Koordinasjon:} Det todelte eierskapet mellom Lobyco og
    Coop Danmark skaper uklarhet om ansvar og prioritering.
\end{itemize}
\keypoint{Modellen minner om Kotters ledelsesmodell, men er mer
analytisk og mindre normativ. Hvor Kotter sier \emph{hva du skal gjøre}
trinn for trinn, beskriver Tawse \& Tabesh \emph{de mekanismene som
avgjør} om implementeringen lykkes, uavhengig av tilnærming. Brukes i
prosjektet til å forsvare \emph{hvorfor} CEO-mandat er nødvendig:
commitment fra øverste ledelse er en av de tre conditions modellen
identifiserer.}
